\section{无监督机器翻译：几何直觉只是启动器}

Lample 的早期研究问题非常苛刻：假设只有两种语言各自的单语文本，没有句子级平行语料，能否学习翻译？课堂口述用一个直观说法描述关键发现：不同语言的词向量空间可能具有相似几何结构，经过变换后能够对齐。这个直觉很重要，但若只写成“两个空间相差一个 rotation”，就会遗漏真正让系统工作的训练闭环。

\subsection{从表示空间到双语词典}

设源语言词向量矩阵为 $X\in\mathbb{R}^{d\times n}$，目标语言锚点矩阵为 $Y\in\mathbb{R}^{d\times n}$。一个经典对齐步骤可写成正交 Procrustes 问题：
\[
W^{\star}=\arg\min_{W}\|WX-Y\|_{F}^{2},
\qquad \text{s.t.}\quad W^{\top}W=I.
\]
其中 $W$ 是保持长度和夹角的正交变换，$\|\cdot\|_{F}$ 是 Frobenius norm。约束 $W^{\top}W=I$ 的直觉是：如果两种语言的局部语义几何相似，就应尽量通过旋转或反射完成对齐，而不是任意拉伸空间。真实系统仍要先找到近似锚点、处理词频和多义性，并面对两种语言结构并不完全同构的问题。

下面的图把“几何启动器”和“完整学习算法”分开。先用单语表示推断粗糙词典，再通过 denoising 与 back-translation 反复改进编码器和解码器。

\lecturefigure{02-cross-lingual-alignment.png}{无监督翻译从表示对齐启动，再通过循环训练逐步提高。}{本地字幕 03:00--05:00；Lample et al., \emph{Unsupervised Machine Translation Using Monolingual Corpora Only}, 2018。}

\begin{knowledgebox}{读图：rotation 为什么不能单独解决翻译}
对齐后的词向量最多提供词级候选，不能自动处理语序、形态变化、长距离依赖和上下文歧义。论文中的系统还需要 denoising auto-encoding 让模型学会语言内部结构，并用 back-translation 把当前译文重新翻回原语言形成自监督信号。几何相似性回答“从哪里开始”，循环训练回答“如何持续变好”。
\end{knowledgebox}

\begin{warningbox}{不要把表示相似误解为语言完全同构}
词向量空间的近似可对齐性依赖语料、词频、训练方法与语言距离。对资源稀缺语言、强形态语言或领域差异很大的语料，粗糙词典可能系统性偏斜。这个方法证明的是“可以构造无监督启动信号”，不是“任何两种语言都只差一个矩阵”。
\end{warningbox}

\teachervoice{老师在这里保留了一个很有价值的研究习惯：先寻找高维问题中的低维结构。词向量看似难以直接观察，但如果假设两个空间共享几何关系，就能把无标注问题转化为可优化的对齐问题。随后再用训练闭环修正这个过强假设。}

\subsection{本章小结}

无监督机器翻译说明，好的研究系统往往由“可解释的启动假设”和“能自我修正的反馈闭环”共同构成。只记住 rotation 会高估假设，忽略 back-translation 又会低估系统工程。

